\section{A contraction for the coefficientwise high-field-degree background}
\label{cdh-section}

The following result concerns only the inherited background
$E(a_p(w))$. It uses the actual pre-gauge representative and the
same incoming and outgoing activity exponent and field radius.
It does not identify a gauge-invariant marginal extraction.

Fix $a_0>0$ and a primitive odd $\ell$. Old carriers $H$ are
connected sets of current $J$-cells; output carriers are
$J$-cells of the next current lattice, or current $JR$-cells,
where $R=\ell^s$ and $J=\ell^j$. The projection has fiber $R^d$
and satisfies
\begin{equation}
 |\pi_RH|\leq2^d+\frac{2^d}{k_R}|H|,
 \qquad k_R=1+\left\lfloor\frac{R-3}{4}\right\rfloor.
 \label{cdh-compression}
\end{equation}
Use compatible finite regulators. Write
\begin{equation}
 \|E\|_{a_0,b}
 =\sup_i\sum_{H\ni i}e^{a_0|H|}\|E_H\|_b,
 \label{cdh-input-norm}
\end{equation}
where each coefficient is holomorphic on its scalar-coordinate
radius-$b$ product polydisc. Each old cell has
$D(J)=dJ^d\dim\mathfrak g$ scalar coordinates.

\subsection{Uniform transport on a larger field-scaling disk}
\label{cdh-scaling-section}

Choose $\kappa\geq a_0+1$ and a finite classical reserve
$A\geq d\kappa$ before choosing $R$. The established pre-gauge
theorem gives constants $C_A,\kappa_0,r_A>0$, fixed at primitive
$\ell$ independently of $R,M$ and volume, with
\begin{equation}
 a_{p,i}(w)=\sum_Xa_{i;X}(w_X),\qquad a_{i;X}(0)=0,\qquad
 \sup_i\sum_Xe^{\kappa_0\tau(\{i\}\cup X)+A|X|}
                   \|a_{i;X}\|_r\leq C_A r/R
 \quad(0<r\leq r_A).
 \label{cdh-ap-input}
\end{equation}
The complete current-field maximum in one source $R$-block is
inside the support sum. Put
\begin{equation}
 L_*=\frac{R}{8C_A}\geq2,\qquad
 0<b\leq\frac{r_A}{L_*},\qquad
 f_\lambda(w)=a_p(\lambda w),\quad |\lambda|<L_*.
 \label{cdh-scaling}
\end{equation}
The absolute rooted norm of $f_\lambda$ on the $b$-ball is at
most $b/8$. Every support truncation therefore has that same
component bound without cancellation.

Cut the rooted supports at distance $\ell_{\rm tail}$, using
exponent $\kappa_0/2$ for the tail, and group roots by $J$.
Provided $2\kappa/J\leq\kappa_0/2$, its macro norm obeys
\begin{equation}
 s_{\rm tail}\leq J^d(b/8)e^{-\kappa_0\ell_{\rm tail}/2}.
 \label{cdh-tail}
\end{equation}
The factor $J^d$ is the sum over old output roots before the
macro-fiber maximum. Restriction to an old $J$-current input
fiber introduces no additional coordinate sum.
For $\ell_{\rm tail}\leq J$ use the conservative connected core
$K_0(H)=N_{\infty,2}(\pi_RH)$. Its constants are
\begin{equation}
 c_0=5^d2^d,\qquad c_1=\frac{5^d2^d}{k_R},\qquad
 h_0=2d,\qquad F=R^d,\qquad m=b/2.
 \label{cdh-core-constants}
\end{equation}
Choose $R$ so that $c_1a_0\leq a_0/2$.
Choose $J$ large and
$\ell_{\rm tail}=\lceil c\log J+C\rceil$ with
$c>4d/\kappa_0$, so that $\ell_{\rm tail}\leq J$ and
\begin{equation}
 t:=s_{\rm tail}/m\leq\tfrac12,\qquad
 D(J)\log\frac1{1-t}\leq a_0/2.
 \label{cdh-tail-budget}
\end{equation}
This is possible because
$D(J)t\leq(d\dim\mathfrak g/4)J^{2d}
e^{-\kappa_0\ell_{\rm tail}/2}\to0$.
Larger $J$ required by other fixed-step modules can be used.
These conditions do not depend on $b$.

The joint core/tail theorem, specialized to no fluctuation
variables and target exponent $a_0$, now gives a canonical
connected family $U_\lambda$ with
\begin{equation}
 \sum_KU_{K,\lambda}(w_K)
   =\sum_HE_H(a_p(\lambda w)|_H),\qquad
 \|U_\lambda\|_{a_0,b}\leq C_0R^d\|E\|_{a_0,b},
 \quad
 C_0=e^{a_0 5^d2^d+2d}S_1.
 \label{cdh-uniform-transport}
\end{equation}
Here $S_1=\sum_{z\in\mathbb Z^d}e^{-|z|_1}$ also bounds the
periodic sum. Support choices are fixed independently of
$\lambda$. The family is holomorphic in the anchored activity
Banach space. The prefactor $C_0$ is independent of $J,R,M$ and
volume; $R^d$ is explicit. Although the background substitution
is nonlinear, this canonical construction is linear in $E$.

\subsection{Coefficientwise fifth-order vanishing and contraction}
\label{cdh-contraction-section}

\begin{theorem}[High-field-degree background transport]
\label{cdh-contraction}
Assume, for every input carrier $H$, the full total field
Taylor tensors satisfy
\begin{equation}
 D^nE_H(0)=0,\qquad n=0,1,2,3,4.
 \label{cdh-jet-condition}
\end{equation}
This includes mixed derivatives in all coordinates of $H$.
Under the preceding parameter choices, the connected output
$\mathcal T_{R,J}E:=U_1$ has the same coefficientwise vanishing
jets and obeys
\begin{equation}
 \|\mathcal T_{R,J}E\|_{a_0,b}
 \leq C_0(8C_A)^5R^{d-5}\|E\|_{a_0,b}.
 \label{cdh-general-contraction}
\end{equation}
In $d=4$, for any prescribed $0<q_{\rm contr}<1$, taking the
finite power $R=\ell^s$ sufficiently large that
\begin{equation}
 R\geq C_0(8C_A)^5/q_{\rm contr}
 \label{cdh-R-contraction}
\end{equation}
in addition to the earlier conditions makes this a genuine
$q_{\rm contr}$ contraction. It is uniform in old depth and
compatible volume.
\end{theorem}

\begin{proof}
In the canonical expansion, a term of tail Taylor degree
$q=|\nu|$ contains
$(\partial^\nu E_H)(c_\lambda)/\nu!$ and $q$ supported tail
factors. Core and tail both vanish at $\lambda=0$ termwise.
For $q\leq4$, condition \eqref{cdh-jet-condition} makes the
derivative at the core vanish to order at least $5-q$ in
$\lambda$, and the tail factors supply order $q$.
For $q\geq5$ the tail factors alone supply order at least five.
This includes the core term $q=0$.
Normal convergence on smaller disks permits differentiation
and grouping: every canonical carrier, and hence $U_\lambda$
as a Banach-valued function, has derivatives of orders zero
through four equal to zero. No cancellation between carriers
is used. The same argument in the field variable $w$ proves
the claimed output jet condition.

For completeness, if a Banach-valued holomorphic $F$ on
$|\lambda|<L$ is bounded by $M$ and has a zero of order $p$ at
zero, then
\begin{equation}
 \|F(\lambda)\|\leq M(|\lambda|/L)^p.
 \label{cdh-Banach-Schwarz}
\end{equation}
The singularity of $F(\lambda)/\lambda^p$ is removable.
On a circle of radius $\rho<L$ its norm is bounded by
$M/\rho^p$; apply the scalar maximum principle to every
bounded linear functional and send $\rho\uparrow L$.
Applying this with $p=5$ to
\eqref{cdh-uniform-transport}, at $\lambda=1$, gives
$C_0R^dL_*^{-5}\|E\|$, which is
\eqref{cdh-general-contraction}. A coefficientwise Cauchy
remainder bound would contain an unnecessary
$(1-L_*^{-1})^{-1}$ factor.
\end{proof}

The parameter order is noncircular: choose the finite reserve
and its constants, then $R$, then $J$ and its tail truncation,
and finally the single radius $b\leq8C_A r_A/R$ together with
any other fixed classical restrictions. The same $b$ is used
at the input and output. For the conservative halo in $d=4$
and $a_0=1/16$, $C_0=e^{633}S_1$; no practically sized scale
is asserted. This result concerns the background transport,
not the new quantum correction or the inherited fluctuation
increment.

\subsection{A bounded Taylor subtraction, with its limitations}
\label{cdh-projection-section}

On the circular product polydisc, define the total homogeneous
Taylor projection by
\begin{equation}
 (P_nE_H)(a)=\frac1{2\pi}\int_0^{2\pi}
               e^{-int}E_H(e^{it}a)\,dt.
 \label{cdh-phase-projection}
\end{equation}
It is the total-degree-$n$ Taylor polynomial, has carrier
$H$, and satisfies $\|P_nE_H\|_b\leq\|E_H\|_b$.
The phase integral estimates the complete homogeneous
polynomial without summing its coordinate monomials.
Consequently
\begin{equation}
 \mathcal R_{\geq5}E_H
     =E_H-\sum_{n=0}^4P_nE_H,\qquad
 \|\mathcal R_{\geq5}E\|_{a_0,b}\leq6\|E\|_{a_0,b}.
 \label{cdh-projection-bound}
\end{equation}
No dimension factor, radius shrink or cardinal reserve is
needed for this projection. The remainder meets
\eqref{cdh-jet-condition}, and in $d=4$
\begin{equation}
 \|\mathcal T_{R,J}\mathcal R_{\geq5}E\|_{a_0,b}
 \leq\frac{6C_0(8C_A)^5}{R}\|E\|_{a_0,b}.
 \label{cdh-projected-bound}
\end{equation}

This is not a gauge-invariant extraction scheme. Local gauge
transformations act nonlinearly and affinely in a link-potential
chart, so the low-degree part and high-degree remainder of a
gauge-invariant total need not be separately gauge invariant.
The rooted/pre-gauge identity applies to the complete
inherited functional before such a split; it cannot be applied
separately to these pieces without another argument.

Nor are the low jets automatically a finite-dimensional space
of couplings. They include kernels on arbitrarily many carrier
geometries and separations. Classifying their local
gauge-invariant content, separating low momentum from field
degree, constructing nonlinear Ward-compatible completions,
and bounding marginal or relevant terms remain separate
requirements. New quantum coefficients may regenerate low
jets. The contraction proved here therefore supplies an
irrelevant-background estimate, not a quantum induction,
running coupling or continuum construction.
